% Báo cáo đồ án IT4863 - Tìm kiếm thông tin
% Biên dịch bằng XeLaTeX (Overleaf: Menu -> Compiler -> XeLaTeX, hoặc dùng latexmkrc đi kèm).
%   xelatex bao-cao.tex && xelatex bao-cao.tex
\documentclass[12pt,a4paper]{report}

\usepackage{fontspec}
\usepackage[vietnamese]{babel}
\usepackage[a4paper,left=3cm,right=2cm,top=2.5cm,bottom=2.5cm]{geometry}
\usepackage{amsmath,amssymb}
\usepackage{graphicx}
\usepackage{booktabs,tabularx,array}
\usepackage{xcolor}
\usepackage{listings}
\usepackage{tikz}
\usetikzlibrary{arrows.meta,positioning,shapes.geometric,fit,calc}
\usepackage{algorithm}
\usepackage{algpseudocode}
\usepackage{enumitem}
\usepackage{caption}
\usepackage{setspace}
\usepackage{fancyhdr}
\usepackage[hidelinks]{hyperref}

\setstretch{1.2}
\emergencystretch=3em
\setlist{nosep,leftmargin=*}
\renewcommand{\arraystretch}{1.25}
\setcounter{tocdepth}{2}

\definecolor{codebg}{RGB}{246,246,246}
\definecolor{codekw}{RGB}{0,70,140}
\definecolor{codecm}{RGB}{110,110,110}
\lstset{
  basicstyle=\ttfamily\footnotesize,
  backgroundcolor=\color{codebg},
  keywordstyle=\color{codekw}\bfseries,
  commentstyle=\color{codecm}\itshape,
  frame=single, rulecolor=\color{black!20},
  breaklines=true, columns=fullflexible, keepspaces=true,
  showstringspaces=false, xleftmargin=4pt, xrightmargin=4pt,
  aboveskip=8pt, belowskip=8pt,
  escapeinside={(*@}{@*)}
}

\newcolumntype{L}{>{\raggedright\arraybackslash}X}

\floatname{algorithm}{Thuật toán}
\renewcommand{\algorithmicrequire}{\textbf{Vào:}}
\renewcommand{\algorithmicensure}{\textbf{Ra:}}

\setlength{\headheight}{15pt}
\pagestyle{fancy}
\fancyhf{}
\fancyhead[L]{\small IT4863 -- Tìm kiếm thông tin}
\fancyhead[R]{\small Nhóm Tìm là thấy}
\fancyfoot[C]{\thepage}

\begin{document}

% ======================================================================
\hypersetup{pageanchor=false}
\begin{titlepage}
\centering
{\large\bfseries ĐẠI HỌC BÁCH KHOA HÀ NỘI}\\[2pt]
{\large\bfseries TRƯỜNG CÔNG NGHỆ THÔNG TIN VÀ TRUYỀN THÔNG}\\[2pt]
\rule{0.5\textwidth}{0.6pt}\\[2.6cm]

{\Large\bfseries BÁO CÁO ĐỒ ÁN MÔN HỌC}\\[8pt]
{\large IT4863 -- Tìm kiếm thông tin}\\[1.6cm]

{\LARGE\bfseries MÁY TÌM KIẾM CHO MIỀN \texttt{HUST.EDU.VN}}\\[10pt]
{\large Crawler theo kiến trúc dịch vụ nhỏ và xếp hạng kết quả trên Lucene}\\[2cm]

\begin{tabular}{@{}ll@{}}
Giảng viên hướng dẫn: & \textbf{TS. Nguyễn Bá Ngọc} \\
Nhóm thực hiện: & \textbf{Tìm là thấy}
\end{tabular}\\[14pt]

\begin{tabular}{lc}
\toprule
Họ và tên & MSHV \\
\midrule
Nguyễn Nam Hán & 20252095M \\
Nguyễn Phúc Nguyên & 20261222M \\
Lê Văn Đức & 20261233M \\
\bottomrule
\end{tabular}

\vfill
{\large Hà Nội, học kỳ 2026.1A}
\end{titlepage}
\hypersetup{pageanchor=true}

% ======================================================================
\pagenumbering{roman}
\chapter*{Tóm tắt}
\addcontentsline{toc}{chapter}{Tóm tắt}

Báo cáo trình bày máy tìm kiếm cho các website thuộc Đại học Bách khoa Hà Nội (\texttt{*.hust.edu.vn}) do nhóm
\emph{Tìm là thấy} xây dựng trong môn IT4863. Hệ thống gồm hai phần dùng chung một cơ sở dữ liệu MongoDB.

\begin{enumerate}
  \item \textbf{Bộ thu thập dữ liệu (crawler)} gồm năm dịch vụ nhỏ trao đổi với nhau qua HTTP/JSON:
        \texttt{mediator} (Python, hàng đợi SQLite, giãn cách theo host, crawl lại định kỳ, phát hiện bài mới qua sitemap và trang danh mục),
        \texttt{downloader} (Node.js, Chromium headless qua Puppeteer, chạy bốn bản),
        \texttt{reader} (Java, jsoup và Apache Tika),
        \texttt{storage} (Spring Boot) và MongoDB. Crawler đã thu thập
        28\,168 trang web từ 95 tên miền con, 4\,690 tài liệu (chủ yếu là PDF), 9\,614 ảnh và 84\,673 cạnh liên kết.
  \item \textbf{Máy tìm kiếm} xây trên Apache Lucene 10. Văn bản tiếng Việt được phân tích thành unigram âm tiết,
        bigram và bản bỏ dấu. Việc xếp hạng chia hai bước: BM25 nhiều trường để lấy ứng viên, sau đó xếp hạng lại theo
        cụm từ và PageRank. Trang rác được phát hiện bằng độ đo ``độ hiếm'' với ngưỡng tự chọn theo thung lũng của
        histogram; trang trùng nội dung được loại bằng mã băm SHA-1. Bốn cấu hình xếp hạng A--D bật dần từng kỹ thuật,
        đi kèm công cụ đánh giá offline bằng pooling, P@10 và nDCG@10.
\end{enumerate}

Cả hệ thống chạy được bằng một lệnh \texttt{docker compose up}. Giao diện web cho phép theo dõi tiến độ crawl, tìm
kiếm và thêm url cần index ngay.

\tableofcontents
\listoffigures
\listoftables

\cleardoublepage
\pagenumbering{arabic}

% ======================================================================
\chapter{Giới thiệu}

\section{Bài toán}

Thông tin của Đại học Bách khoa Hà Nội nằm rải rác trên gần một trăm website: cổng chính \texttt{hust.edu.vn}, website
của các trường, khoa, viện (\texttt{soict}, \texttt{seee}, \texttt{bme}, \texttt{bktextile}\ldots), thư viện, phòng ban
và các trang dự án. Những thông tin sinh viên hay cần như học phí, lịch thi, quy chế đào tạo thường nằm trong các tệp
PDF, DOCX đính kèm chứ không phải trong trang web. Các công cụ tìm kiếm tổng quát trả về lẫn cả website ngoài trường
và ít khi đánh chỉ mục nội dung tệp đính kèm.

Đồ án xây dựng một máy tìm kiếm riêng cho miền \texttt{hust.edu.vn} và làm đủ các khâu của một hệ tìm kiếm web: thu
thập, bóc nội dung, lưu trữ, đánh chỉ mục, xếp hạng và đánh giá.

\section{Mục tiêu}

\begin{itemize}
  \item Tự động thu thập trang web, tài liệu và ảnh thuộc \texttt{*.hust.edu.vn}, mở rộng theo liên kết, không gửi
        yêu cầu dồn dập vào một máy chủ (giới hạn tốc độ theo host) và tự cập nhật định kỳ.
  \item Lấy được phần nội dung chính của trang, bỏ menu và chân trang; đọc được chữ trong PDF, DOCX, XLSX.
  \item Lưu nội dung và đồ thị liên kết vào MongoDB theo một lược đồ thống nhất.
  \item Tìm kiếm tiếng Việt chấp nhận cả truy vấn có dấu và không dấu, có dùng tín hiệu liên kết, lọc trang rác và
        trang trùng.
  \item Đo được mỗi kỹ thuật xếp hạng đóng góp bao nhiêu bằng đánh giá offline.
\end{itemize}

\section{Cấu trúc báo cáo}

Chương~\ref{ch:arch} nói về kiến trúc chung. Chương~\ref{ch:crawler} trình bày bộ thu thập dữ liệu và số liệu thu
được. Chương~\ref{ch:search} trình bày máy tìm kiếm: phân tích văn bản, chỉ mục, BM25, PageRank, lọc rác, loại trang
trùng và truy vấn hai bước. Chương~\ref{ch:eval} nói về cách đánh giá, Chương~\ref{ch:ops} về triển khai và kiểm thử,
Chương~\ref{ch:limit} về hạn chế, hướng phát triển và kết luận.

% ======================================================================
\chapter{Kiến trúc tổng thể}\label{ch:arch}

\section{Sơ đồ hệ thống}

Hệ thống gồm sáu dịch vụ. Mỗi dịch vụ có một thư mục và một \texttt{Dockerfile} riêng, trao đổi với nhau bằng HTTP và
JSON (Hình~\ref{fig:arch}). Dữ liệu đi theo một chiều: \texttt{mediator} $\rightarrow$ \texttt{downloader}
$\rightarrow$ \texttt{reader} $\rightarrow$ \texttt{storage} $\rightarrow$ MongoDB. Vòng lặp crawl khép lại ở chỗ
\texttt{reader} gửi các url \texttt{hust.edu.vn} mới tìm thấy về lại \texttt{mediator}. Dịch vụ \texttt{search} đọc
MongoDB để dựng chỉ mục Lucene và trả lời truy vấn.

\begin{figure}[ht]
\centering
\begin{tikzpicture}[
  svc/.style={draw, rounded corners=3pt, fill=blue!6, align=center, minimum width=2.5cm, minimum height=1.2cm, font=\small},
  db/.style={draw, cylinder, shape border rotate=90, aspect=0.25, fill=yellow!15, align=center, minimum width=2.2cm, minimum height=1.3cm, font=\small},
  ext/.style={draw, dashed, rounded corners=3pt, align=center, font=\small, minimum height=0.9cm},
  arr/.style={-{Stealth[length=2.2mm]}, thick},
  lbl/.style={font=\scriptsize, midway, inner sep=2pt}
]
\node[svc] (med) {mediator\\\scriptsize :5200\\\scriptsize Python, SQLite};
\node[svc, right=1.9cm of med] (dl) {downloader $\times 4$\\\scriptsize :8800\\\scriptsize Node, Puppeteer};
\node[svc, right=1.9cm of dl] (rd) {reader\\\scriptsize :6600\\\scriptsize Java, jsoup, Tika};
\node[svc, below=2.2cm of rd] (st) {storage\\\scriptsize :2222\\\scriptsize Spring Boot};
\node[db, left=1.9cm of st] (mg) {MongoDB\\\scriptsize :27018};
\node[svc, left=1.9cm of mg, fill=green!8] (se) {search\\\scriptsize :7700\\\scriptsize Lucene 10};
\node[ext, above=1.0cm of med] (seed) {\texttt{feed.py}\\\texttt{link\_\allowbreak{}all.txt}};
\node[ext, below=1.0cm of se] (ui) {\texttt{dashboard.html}\\\scriptsize trình duyệt};
\node[ext, above=1.0cm of dl] (web) {Website\\\texttt{*.hust.edu.vn}};

\draw[arr] (seed) -- node[lbl,right] {POST /add (url mầm)} (med);
\draw[arr] (med) -- node[lbl,above] {/ping, /load} (dl);
\draw[arr] (dl) -- node[lbl,above] {POST /parse} (rd);
\draw[arr, <->] (dl) -- node[lbl,right] {HTTP GET} (web);
\draw[arr] (rd) -- node[lbl,right] {POST /add} (st);
\draw[arr] (st) -- node[lbl,above] {upsert} (mg);
\draw[arr] (mg) -- node[lbl,above] {dựng chỉ mục} (se);
\draw[arr] (ui) -- node[lbl,right] {GET /search} (se);
\draw[arr, dashed] ([xshift=-6mm]rd.south) -- ++(0,-1.0) -| node[lbl,above,pos=0.25] {POST /add \{url hust mới\}} (med.south);
\end{tikzpicture}
\caption{Kiến trúc tổng thể của hệ thống}\label{fig:arch}
\end{figure}

\section{Các dịch vụ}

\begin{table}[ht]
\centering
\caption{Các dịch vụ và công nghệ}\label{tab:services}
\begin{tabularx}{\textwidth}{l L c}
\toprule
Dịch vụ (thư mục) & Vai trò, công nghệ & Cổng \\
\midrule
\texttt{mediator} & Giữ hàng đợi url (SQLite), giao url cho các bản downloader theo vòng, giãn cách theo host, hẹn lịch crawl lại, phát hiện bài mới. Python 3.13, \texttt{asyncio}, \texttt{aiohttp}. & 5200 \\
\texttt{downloader} & Dùng Chromium headless (Puppeteer) cho trang web, \texttt{fetch} cho tài liệu và ảnh. Mỗi bản chỉ tải một url mỗi lúc. Node.js 24. & 8800 \\
\texttt{reader} & Lấy tiêu đề, khối nội dung chính, liên kết; đọc chữ trong tệp bằng Apache Tika; chuyển kết quả sang \texttt{storage} và \texttt{mediator}. Java 21. & 6600 \\
\texttt{storage} & Kiểm tra và ghi bản ghi vào \texttt{html}/\texttt{other}/\texttt{picture}, ghi cạnh vào \texttt{graph}. Java 21, Spring Boot 3.3. & 2222 \\
\texttt{search} & Dựng chỉ mục, xếp hạng, API \texttt{/search}, phục vụ giao diện web. Java 21, Lucene 10.5. & 7700 \\
MongoDB & Cơ sở dữ liệu dùng chung, MongoDB 7. & 27018 \\
\bottomrule
\end{tabularx}
\end{table}

Nhóm chia thành nhiều dịch vụ nhỏ vì ba lý do. Thứ nhất, mỗi khâu tốn tài nguyên theo một kiểu: tải trang chủ yếu chờ
mạng và tốn RAM cho trình duyệt, đọc PDF tốn CPU, nên cần nhân bản riêng từng khâu (hệ chạy bốn bản
\texttt{downloader}, các khâu khác một bản). Thứ hai, mỗi khâu dùng được ngôn ngữ có thư viện tốt nhất cho nó:
Puppeteer trên Node.js, jsoup, Tika và Lucene trên Java. Thứ ba, giữa các khâu chỉ có vài endpoint JSON nên kiểm thử
và thay thế từng phần khá dễ.

\section{Mô hình dữ liệu}

Dữ liệu nằm trong database \texttt{HUST\_\allowbreak{}2026\_\allowbreak{}IT4863\_\allowbreak{}TEST} gồm bốn collection (Bảng~\ref{tab:mongo}). Lược đồ được
ràng buộc bằng \texttt{\$jsonSchema} ở mức \texttt{strict} (tệp \path{storage/document/ddl_create_document.js}):
\texttt{url} phải bắt đầu bằng \texttt{http(s)://}, \texttt{base64} phải là chuỗi Base64 hợp lệ. Mỗi collection nội
dung có chỉ mục duy nhất trên \texttt{url}; \texttt{graph} có chỉ mục duy nhất trên cặp
(\texttt{source}, \texttt{destinations}) và thêm một chỉ mục phụ trên \texttt{destinations}.

\begin{table}[ht]
\centering
\caption{Các collection trong MongoDB}\label{tab:mongo}
\begin{tabularx}{\textwidth}{l L}
\toprule
Collection & Nội dung \\
\midrule
\texttt{html} & Trang web: \texttt{url}, \texttt{format}, \texttt{title}, \texttt{content} (văn bản đã bóc), \texttt{createdAt}, \texttt{updatedAt}. \\
\texttt{other} & Tài liệu (pdf, docx, xlsx\ldots): các trường như trên, có thể thêm \texttt{origin}. \\
\texttt{picture} & Ảnh: các trường như trên, thêm \texttt{base64} chứa byte ảnh. \\
\texttt{graph} & Mỗi bản ghi là một cạnh liên kết: \texttt{source}, \texttt{destinations} (một url đích), \texttt{createdAt}, \texttt{updatedAt}. \\
\bottomrule
\end{tabularx}
\end{table}

% ======================================================================
\chapter{Thu thập dữ liệu}\label{ch:crawler}

Chương này đi theo đường đi của một url qua bốn dịch vụ: \texttt{mediator} quyết định url nào được tải và khi nào;
\texttt{downloader} lấy nội dung; \texttt{reader} bóc chữ và liên kết; \texttt{storage} ghi vào MongoDB.

\section{Dịch vụ mediator}

\subsection{Vai trò}

Mỗi bản \texttt{downloader} chỉ tải một url một lúc. Nếu đẩy thẳng hàng nghìn url vào thì nó quá tải, còn nếu gửi dồn
vào một website thì website đó có thể chặn. \texttt{mediator} đứng giữa để giải quyết cả hai: giữ mọi url trong một hàng
đợi bền vững, bỏ url trùng, giãn cách các lần gửi vào cùng một host, chỉ giao url khi có bản \texttt{downloader} rảnh
và hẹn ngày kiểm tra lại từng url.

\texttt{mediator} chạy hai luồng bất đồng bộ trong cùng một tiến trình \texttt{asyncio}, gặp nhau ở hàng đợi:

\begin{enumerate}
  \item \textbf{Luồng nhận} (\texttt{POST /add}): chuẩn hoá, phân loại và cất url vào hàng đợi, trả lời ngay mà không
        chờ tải.
  \item \textbf{Luồng giao}: vòng lặp liên tục, mỗi lượt lấy tối đa một url đến hạn và giao cho một bản
        \texttt{downloader} đang rảnh.
\end{enumerate}

Tách như vậy thì bên gửi url gửi bao nhiêu cũng được trả lời ngay, còn tốc độ giao do tốc độ tải quyết định.

\subsection{Chuẩn hoá url}

Cùng một trang có thể viết bằng nhiều cách. Hàm \texttt{canonicalize\_\allowbreak{}url} đưa chúng về một dạng duy nhất để không tải
một trang hai lần:

\begin{itemize}
  \item viết thường scheme và tên miền, bỏ dấu chấm cuối tên miền, bỏ cổng mặc định (80 với http, 443 với https);
  \item chuẩn hoá đường dẫn (bỏ \texttt{.} và \texttt{..}), giữ dấu \texttt{/} cuối nếu có;
  \item bỏ phần \texttt{\#fragment};
  \item bỏ các tham số theo dõi và phiên: \texttt{utm\_\allowbreak{}*}, \texttt{fbclid}, \texttt{gclid}, \texttt{sessionid},
        \texttt{jsessionid}, \texttt{phpsessid};
  \item sắp xếp các tham số còn lại theo khoá rồi theo giá trị.
\end{itemize}

\begin{sloppypar}
Ví dụ, \texttt{HTTPS://Hust.edu.vn:443/a/../b?z=2\&a=1\&\allowbreak utm\_\allowbreak{}source=fb\#top} trở thành
\texttt{https://hust.edu.vn/b?\allowbreak a=1\&z=2}. Url bị loại nếu không phải \texttt{http}/\texttt{https}, không có host, có
thông tin đăng nhập (\texttt{user:pass@}) hoặc có hơn 50 cặp tham số (thường là bẫy crawler sinh url vô hạn).
\end{sloppypar}

\subsection{Phân loại url và chu kỳ crawl lại}

Mỗi url thuộc một trong ba loại vì chúng thay đổi với tốc độ khác nhau (Bảng~\ref{tab:kinds}). Các chu kỳ chỉnh được
trong tệp cấu hình YAML; giá trị 0 nghĩa là không crawl lại.

\begin{table}[ht]
\centering
\caption{Phân loại url và chu kỳ kiểm tra lại}\label{tab:kinds}
\begin{tabularx}{\textwidth}{l L c}
\toprule
Loại & Cách nhận biết & Kiểm tra lại sau \\
\midrule
\texttt{hub} (trang chủ) & Đường dẫn rỗng hoặc \texttt{/}; tin mới đăng liên tục. & 1 ngày \\
\texttt{page} (trang thường) & Các url còn lại. & 7 ngày \\
\texttt{document} (tài liệu, ảnh) & Đuôi \texttt{.pdf}, \texttt{.doc(x)}, \texttt{.xls(x)}, \texttt{.ppt(x)}, \texttt{.odt}, \texttt{.zip}, \texttt{.jpg}, \texttt{.png}\ldots; hiếm khi đổi. & 30 ngày \\
\bottomrule
\end{tabularx}
\end{table}

\subsection{Hàng đợi bền vững trên SQLite}

Hàng đợi là một tệp SQLite có hai bảng. Bảng \texttt{urls} lưu mọi url đã biết cùng loại, hạn giao kế tiếp
(\texttt{next\_\allowbreak{}due\_\allowbreak{}at}), lần gửi cuối và số lần lỗi liên tiếp. Bảng \texttt{hosts} lưu thời điểm gửi gần nhất vào từng
host.

\begin{lstlisting}[language=SQL, caption={Lược đồ hàng đợi}]
CREATE TABLE urls (
    url TEXT PRIMARY KEY, host TEXT NOT NULL, kind TEXT NOT NULL,
    next_due_at REAL, last_sent_at REAL, attempts INTEGER NOT NULL DEFAULT 0);
CREATE INDEX urls_next_due ON urls (next_due_at);
CREATE TABLE hosts (host TEXT PRIMARY KEY, last_sent_at REAL NOT NULL);
\end{lstlisting}

Thêm url dùng \texttt{INSERT OR IGNORE}: url đã có thì giữ nguyên lịch cũ, nên gửi lại url cũ không làm nó bị tải thêm.
Mọi thao tác ghi đều commit ngay, vì vậy tắt dịch vụ rồi bật lại thì việc giao url tiếp tục đúng chỗ đã dừng. Việc chọn
url kế tiếp là một truy vấn duy nhất, xét cả điều kiện đến hạn lẫn khoảng nghỉ của host:

\begin{lstlisting}[language=SQL, caption={Chọn url kế tiếp}]
SELECT u.* FROM urls u LEFT JOIN hosts h ON h.host = u.host
WHERE u.next_due_at <= :now
  AND (h.last_sent_at IS NULL OR h.last_sent_at + :host_delay <= :now)
ORDER BY u.next_due_at <= 0 DESC,
         u.kind = 'document' DESC,
         u.next_due_at
LIMIT 1;
\end{lstlisting}

Thứ tự ưu tiên trong \texttt{ORDER BY} là: url được đánh dấu ưu tiên (\texttt{next\_\allowbreak{}due\_\allowbreak{}at = 0}), rồi tài liệu và ảnh, rồi url đến hạn sớm nhất. Khoảng nghỉ nằm ngay trong mệnh đề \texttt{WHERE}: host chưa nghỉ đủ \texttt{host\_\allowbreak{}delay\_\allowbreak{}seconds} (cấu hình 2 giây) bị
bỏ qua và truy vấn lấy url của host khác. Nhờ vậy khi một host đang nghỉ, các host khác vẫn được thu thập xen kẽ, và
không host nào nhận quá một yêu cầu mỗi hai giây từ toàn hệ thống. Tài liệu được ưu tiên trước trang thường vì đó là
phần thông tin có giá trị nhất (quy chế, thông báo) và tải bằng \texttt{fetch} nhanh hơn dựng trang bằng trình duyệt.

\subsection{Luồng giao url}

\texttt{mediator} làm việc với \texttt{downloader} theo hai bước: hỏi \texttt{GET /ping} xem đã rảnh chưa (trả
\texttt{\{"status":"đang chờ"\}} hoặc \texttt{"đang bận"}), nếu rảnh thì gửi \texttt{POST /load \{"url": ...\}}. Các bản
\texttt{downloader} được duyệt theo vòng (round robin), bắt đầu từ bản sau bản vừa dùng và bỏ qua bản đang bị cho nghỉ
vì lỗi trước đó. Thuật toán~\ref{alg:dispatch} mô tả một lượt giao.

\begin{algorithm}[ht]
\caption{Một lượt giao url của mediator}\label{alg:dispatch}
\begin{algorithmic}[1]
\State $u \gets$ url đến hạn kế tiếp (truy vấn ở trên)
\If{$u = \varnothing$} \State \Return \textsc{false} \Comment{nghỉ 0,5 giây rồi thử lại}
\EndIf
\For{mỗi node $n$ theo thứ tự vòng, bỏ node đang nghỉ}
  \State $s \gets$ \textsc{Ping}$(n)$ \Comment{hết hạn 2 giây; chuẩn hoá NFC}
  \If{không có phản hồi} \State cho $n$ nghỉ 30 giây; \textbf{continue}
  \ElsIf{$s \neq$ ``đang chờ''} \State \textbf{continue}
  \EndIf
  \State $r \gets$ \textsc{Load}$(n, u)$ \Comment{hết hạn 10 giây}
  \If{$r = 409$} \State \textbf{continue} \Comment{node vừa nhận việc khác}
  \ElsIf{$r = 2xx$}
    \State $u.\mathit{next} \gets now + \mathit{recrawl}(u.\mathit{kind})$; ghi thời điểm gửi của host
    \State \Return \textsc{true}
  \ElsIf{$r = 4xx$} \Comment{lỗi do url, node vẫn khoẻ}
    \State $u.\mathit{next} \gets now + \mathit{recrawl}(u.\mathit{kind})$; \Return \textsc{false}
  \Else \Comment{5xx hoặc hết giờ}
    \State cho $n$ nghỉ 30 giây; $u.\mathit{attempts} \gets u.\mathit{attempts} + 1$
    \If{$u.\mathit{attempts} \ge 5$} \State $u.\mathit{next} \gets now + \mathit{recrawl}(u.\mathit{kind})$; $u.\mathit{attempts} \gets 0$
    \Else \State $u.\mathit{next} \gets now$ \Comment{thử lại ở lượt sau}
    \EndIf
    \State \Return \textsc{false}
  \EndIf
\EndFor
\State \Return \textsc{false}
\end{algorithmic}
\end{algorithm}

Một số chi tiết trong cách xử lý:

\begin{itemize}
  \item \textbf{Chuẩn hoá Unicode của trạng thái.} Chuỗi ``đang chờ'' có thể đến ở dạng NFD (chữ gốc cộng dấu rời) nếu
        \texttt{downloader} được viết trên hệ điều hành khác, nên \texttt{mediator} đưa về NFC trước khi so sánh.
  \item \textbf{Cho node lỗi nghỉ.} Node không trả lời \texttt{/ping} trong 2 giây, hoặc lỗi khi \texttt{/load}, bị
        bỏ qua 30 giây để không gọi liên tục vào một dịch vụ đang hỏng. Node sống lại thì việc giao tự tiếp tục.
  \item \textbf{Phân biệt lỗi của url và lỗi của node.} Mã 4xx nghĩa là node vẫn khoẻ nhưng từ chối url đó, nên node
        không bị cho nghỉ còn url được lùi một chu kỳ thay vì thử lại ngay.
  \item \textbf{Giới hạn số lần thử.} Url lỗi 5 lần liên tiếp bị lùi một chu kỳ, nên url hỏng được thử vài lần chứ
        không lặp mãi.
  \item \textbf{Ưu tiên.} \texttt{POST /add} kèm \texttt{"priority": true} đặt \texttt{next\_\allowbreak{}due\_\allowbreak{}at = 0} cho các url
        gửi lên (kể cả url đã biết) để chúng được tải ngay. Chức năng ``Thêm url'' trên giao diện dùng cơ chế này.
\end{itemize}

\subsection{Duy trì dữ liệu và tải bài mới}\label{sec:incremental}

Chu kỳ crawl lại ở Bảng~\ref{tab:kinds} chỉ cho biết khi nào tải lại một url đã biết. Còn bài viết mới thì chưa có url
nào trong hàng đợi để mà tải lại, nên cần hai cơ chế bổ sung để phát hiện chúng.

\paragraph{Đọc sitemap.} Định kỳ \texttt{mediator} lấy tệp \texttt{sitemap.xml} của từng tên miền con. Url chưa có
trong SQLite được thêm vào hàng đợi để tải ngay. Url đã có nhưng \texttt{<lastmod>} mới hơn lần tải trước thì bị đặt lại
\texttt{next\_due\_at = 0} để tải lại sớm. Không phải website nào cũng có sitemap đầy đủ, nên cơ chế này không thay thế
được cơ chế thứ hai bên dưới.

\paragraph{Duyệt trang danh mục theo phân trang.} Website khoa, viện thường xếp tin tức thành danh mục nhiều trang
(\texttt{/tin-tuc?page=2}). Khi có bài mới, chúng xuất hiện ở trang 1 và đẩy bài cũ xuống các trang sau. Hằng ngày
\texttt{mediator} tải lại trang 1 (loại \texttt{hub}, Bảng~\ref{tab:kinds}). Nếu toàn bộ bài trên trang 1 đều mới thì có
thể còn bài mới ở trang 2, nên cần tải tiếp; nếu đã gặp bài cũ thì dừng. Thuật toán~\ref{alg:incremental} mô tả cách này.

\begin{algorithm}[ht]
\caption{Tải bài mới theo trang danh mục}\label{alg:incremental}
\begin{algorithmic}[1]
\Require trang danh mục đầu tiên $u_1$; số trang tối đa $K_{\max} = 10$
\Ensure mọi bài mới được đưa vào hàng đợi; dừng khi gặp bài đã biết
\State $k \gets 1$
\While{$k \le K_{\max}$}
  \State tải và bóc trang $u_k$; $L_k \gets$ các url bài viết trong khối nội dung chính
  \State $m \gets$ số url trong $L_k$ chưa có trong bảng \texttt{urls}
  \State đưa các url mới đó vào hàng đợi để tải ngay
  \If{$m < |L_k|$} \State \textbf{break} \Comment{đã gặp bài cũ: các bài mới nằm hết ở các trang đã duyệt}
  \EndIf
  \State $u_{k+1} \gets$ url của trang kế tiếp (tăng tham số trang lên $k+1$)
  \State đặt $u_{k+1}$ ở mức ưu tiên; $k \gets k + 1$
\EndWhile
\end{algorithmic}
\end{algorithm}

Có ba trường hợp tại trang $k$. Nếu $m = 0$, không có bài mới và việc duyệt dừng ngay ở trang 1. Nếu $0 < m < |L_k|$,
trang chứa cả bài mới lẫn bài cũ, nghĩa là các bài mới đã nằm trọn trong những trang vừa duyệt nên dừng. Chỉ khi
$m = |L_k|$, tức cả trang đều mới, mới có khả năng còn bài mới ở trang sau và thuật toán sang trang $k+1$. Giới hạn
$K_{\max}$ chặn trường hợp một website đăng cả loạt bài cùng lúc, hoặc đổi cách đánh số trang khiến mọi trang đều
được coi là mới. Nhờ điều kiện dừng này, mỗi lần cập nhật thường chỉ tải một trang danh mục thay vì toàn bộ chuỗi.

\subsection{API và cấu trúc mã}

\texttt{mediator} có bốn endpoint: \texttt{POST /add} (nhận một gói url, trả số url mới, trùng, bị loại và url đã
chuẩn hoá của từng url), \texttt{POST /stop} và \texttt{POST /start} (tạm dừng hoặc chạy tiếp việc giao, hàng đợi giữ
nguyên), \texttt{GET /status} (trạng thái, tổng số url, số url đến hạn). Cứ 30 giây dịch vụ in một dòng tóm tắt vào log.

\begin{lstlisting}[caption={Ví dụ gọi POST /add}]
$ curl -X POST localhost:5200/add -d '{"urls": ["https://hust.edu.vn/",
       "https://HUST.edu.vn:443/#top", "mailto:a@hust.edu.vn"]}'
{"accepted": 1, "duplicates": 1, "rejected": 1,
 "urls": ["https://hust.edu.vn/", "https://hust.edu.vn/", null]}
\end{lstlisting}

Mã nguồn chia lớp theo kiểu ports-and-adapters:

\begin{itemize}
  \item \texttt{domain/}: logic thuần, không I/O: chuẩn hoá url, phân loại, lập lịch lại, xoay vòng node.
  \item \texttt{application/}: use case \texttt{add\_\allowbreak{}urls} và \texttt{Dispatcher}, chỉ phụ thuộc vào các giao diện
        (port) \texttt{UrlQueue}, \texttt{Downloader}, \texttt{Clock}.
  \item \texttt{infrastructure/}: cài đặt các port bằng SQLite và \texttt{aiohttp}; đọc cấu hình YAML.
  \item \texttt{api/}: HTTP server \texttt{aiohttp}.
\end{itemize}

Chia như vậy cho phép kiểm thử luồng giao bằng đồng hồ giả và \texttt{downloader} giả, không cần mạng hay chờ thời gian
thật. Một bài kiểm thử kiến trúc (\texttt{test\_\allowbreak{}architecture.py}) bảo đảm lớp \texttt{domain} không import lớp ngoài.

\section{Dịch vụ downloader}

\texttt{downloader} là một HTTP server Node.js. Nó nhận url qua \texttt{POST /load}, trả \texttt{202} ngay rồi tải ở
nền, xong thì gửi kết quả thẳng sang \texttt{reader} (\texttt{POST /parse}). Mỗi bản chỉ tải một url một lúc: khi đang
tải, \texttt{/ping} trả ``đang bận'' và \texttt{/load} trả \texttt{409}. Hệ thống chạy bốn bản
(\texttt{deploy.replicas: 4}), địa chỉ của chúng khai báo trong cấu hình của \texttt{mediator}.

\subsection{Tải trang web bằng trình duyệt headless}

Nhiều website của trường nạp nội dung bằng JavaScript nên tải HTML thô là không đủ. \texttt{downloader} dùng Puppeteer~\cite{puppeteer}
điều khiển một tiến trình Chromium dùng chung, mỗi lần tải mở một tab:

\begin{enumerate}
  \item Bật chặn request, huỷ mọi yêu cầu loại \texttt{image}, \texttt{media}, \texttt{font}, \texttt{stylesheet}.
        \texttt{reader} chỉ cần DOM, nên việc này vừa nhanh hơn vừa bớt request vào website.
  \item Gọi \texttt{page.goto(url, \{waitUntil: 'load'\})} với hạn 30 giây, sau đó chờ mạng rảnh 300 ms (thêm tối đa
        2 giây) để các lời gọi AJAX nạp xong nội dung.
  \item Nếu mã HTTP $\ge 400$ thì ghi log rồi bỏ. Nếu \texttt{Content-Type} không phải \texttt{text/html} (ví dụ url
        không có đuôi nhưng trả về PDF), hoặc Chromium huỷ điều hướng (\texttt{ERR\_\allowbreak{}ABORTED}) vì đây là tệp tải xuống,
        thì chuyển sang tải nhị phân.
  \item Gửi \texttt{page.content()} (DOM sau khi chạy JavaScript) cùng \texttt{final\_\allowbreak{}url} sau chuyển hướng sang
        \texttt{reader}.
\end{enumerate}

Khi Chromium chết hoặc mất kết nối, cache của trình duyệt bị xoá để lần tải sau khởi động lại sạch.

\subsection{Tải tài liệu và ảnh}

Url có đuôi tài liệu hoặc ảnh (\texttt{pdf}, \texttt{docx}, \texttt{xlsx}, \texttt{pptx}, \texttt{jpg},
\texttt{png}\ldots) được tải thẳng bằng \texttt{fetch}, theo chuyển hướng, giới hạn 50 MB (kiểm tra cả
\texttt{Content-Length} lẫn số byte thực nhận). Nếu có \texttt{Content-Type}, byte của tệp được gửi nguyên trong body
\texttt{/parse?url=...\&final\_\allowbreak{}url=...}, tránh việc Base64 làm dữ liệu to thêm 33\%; nếu không thì gửi gói JSON với
\texttt{encoding: "base64"}.

\section{Dịch vụ reader}

\texttt{reader} nhận trang hoặc tệp, bóc chữ và liên kết rồi chuyển tiếp: bản ghi sang \texttt{storage}, url
\texttt{hust.edu.vn} mới sang \texttt{mediator}. Dịch vụ viết bằng Java 21, dùng HTTP server có sẵn của JDK, jsoup để
phân tích HTML và Apache Tika 4~\cite{tika} để đọc tệp.

\subsection{Lấy tiêu đề}

Tiêu đề lấy theo thứ tự ưu tiên: \texttt{<meta property="og:title">}, thẻ \texttt{<h1>} đầu tiên, rồi \texttt{<title>}.
Thẻ \texttt{<title>} thường kèm tên website (``Học phí -- Trường CNTT\&TT'') nên để sau cùng. Nếu không có tiêu đề nào
thì lấy tên tệp trong url.

\subsection{Tìm khối nội dung chính}

Trang của trường có nhiều phần lặp lại như menu, thanh bên, chân trang. Nếu index cả trang thì trang nào cũng khớp với
các từ trong menu. \texttt{reader} chọn \emph{khối nội dung chính} bằng một heuristic chấm điểm theo mật độ chữ, cùng
ý tưởng với các thuật toán kiểu Readability:

\begin{enumerate}
  \item Xoá các phần không hiển thị: \texttt{script}, \texttt{style}, \texttt{noscript}, \texttt{template},
        \texttt{iframe} và phần tử có thuộc tính \texttt{hidden}.
  \item Với mỗi phần tử $e$ có chữ trực tiếp (\texttt{ownText}, không tính chữ trong thẻ \texttt{<a>}) dài
        $n \ge 25$ ký tự: leo lên thẻ khối gần nhất $e'$ (một \texttt{<span>} nằm trong \texttt{<p>} thì tính từ
        \texttt{<p>}), rồi cộng điểm
        \[
          \mathit{score}(\mathrm{cha}(e')) \mathrel{+}= n, \qquad
          \mathit{score}(\mathrm{ông}(e')) \mathrel{+}= n/2 .
        \]
  \item Khối nội dung là phần tử có điểm cao nhất; nếu không có đoạn chữ nào thì lấy \texttt{<body>}.
\end{enumerate}

Đoạn chữ ngắn hơn 25 ký tự (mục menu, nhãn nút) không được tính, và chữ nằm trong liên kết cũng vậy, nên một menu toàn
liên kết không thể thắng phần thân bài viết. Cộng thêm một nửa điểm cho phần tử ông giúp chọn được vùng bao ngoài khi
bài viết bị chia thành nhiều \texttt{<div>} con.

\subsection{Liên kết và văn bản liên kết}

\texttt{reader} tách liên kết thành hai tập:

\begin{itemize}
  \item \textbf{Liên kết nằm trong khối nội dung} (\texttt{<a href>} và \texttt{<img src>}) được ghi vào đồ thị
        \texttt{graph}. Văn bản liên kết là \emph{câu chứa liên kết} chứ không chỉ chữ trong thẻ: \texttt{reader} duyệt
        các nút chữ của khối câu gần nhất (\texttt{p}, \texttt{li}, \texttt{td}, \texttt{h1}--\texttt{h6}\ldots) và cắt
        từ dấu kết câu gần nhất bên trái đến dấu kết câu đầu tiên bên phải. Nhờ vậy liên kết ``xem tại đây'' vẫn mang
        ngữ cảnh ``Quy chế đào tạo 2025, xem tại đây''. Với ảnh, văn bản lấy từ \texttt{alt}, không có thì lấy
        \texttt{title}.
  \item \textbf{Mọi liên kết của trang}, kể cả menu và chân trang, được lọc lấy các url có host \texttt{hust.edu.vn}
        hoặc \texttt{*.hust.edu.vn} rồi gửi về \texttt{mediator}. Menu cũng được dùng để mở rộng crawl, vì menu chính
        là bản đồ của website.
\end{itemize}

Liên kết được đổi thành url tuyệt đối theo \texttt{final\_\allowbreak{}url}, bỏ fragment, bỏ liên kết về chính trang và bỏ trùng. Ký
tự không hợp lệ (dấu cách, chữ có dấu) được mã hoá \texttt{\%XX} để \texttt{storage} không từ chối cả bản ghi.

\subsection{Đọc tài liệu và ảnh}

Tệp nhị phân được đưa qua \texttt{AutoDetectParser} của Tika, nó tự nhận dạng định dạng theo nội dung chứ không tin vào
đuôi tệp. Kết quả:

\begin{itemize}
  \item \textbf{Tài liệu} (PDF, DOC/DOCX, XLS/XLSX, PPT/PPTX): lấy chữ thuần, cắt ở 1 triệu ký tự. Tiêu đề lấy từ
        metadata của tệp, không có thì lấy tên tệp. Định dạng ghi theo đuôi tệp, không có đuôi thì suy từ MIME.
  \item \textbf{Ảnh}: không OCR. Hệ thống lưu tiêu đề (tên tệp), metadata đã chuẩn hoá (MIME, kích thước điểm ảnh, ngày
        chụp, tác giả, máy ảnh) và byte ảnh dạng Base64. Ảnh trên 11 MB bị bỏ qua, vì Base64 làm dữ liệu to thêm một
        phần ba, cộng với các trường khác sẽ vượt giới hạn 16 MB của một bản ghi MongoDB.
\end{itemize}

Tệp không đọc được trả \texttt{415} hoặc \texttt{422}; trang không bóc được chữ trả \texttt{422}. Bảng mã của trang lấy
từ \texttt{charset} trong \texttt{Content-Type}, mặc định UTF-8.

\subsection{Chuyển tiếp}

Bản ghi được gửi bằng \texttt{POST /add} sang \texttt{storage}, danh sách url hust gửi sang \texttt{POST /add} của
\texttt{mediator}. Mỗi lời gọi thử lại một lần với hạn 10 giây; nếu vẫn thất bại thì \texttt{/parse} trả \texttt{502}
kèm tên dịch vụ và lý do. Mỗi request ghi một dòng log gồm url, định dạng, mã trạng thái, thời gian xử lý và số liên
kết.

\section{Dịch vụ storage}

\texttt{storage} là ứng dụng Spring Boot với một endpoint ghi (\texttt{POST /add}) và hai endpoint đọc:
\texttt{GET /stats} (số bản ghi, dung lượng từng collection, bản ghi mới nhất) và \texttt{GET /record} (xem một bản
ghi theo url).

\begin{lstlisting}[caption={Body của POST /add trên storage}]
{"url": "https://hust.edu.vn/tin-tuc1", "format": "html", "title": "Tin (*@tức@*) 1",
 "content": "(*@Cơ@*) (*@hội@*) (*@thực@*) (*@tập@*) (*@có@*) (*@lương@*) (*@tại@*) (*@Nhật@*) cho sinh (*@viên@*) HUST.",
 "out-links": [{"url": "https://hust.edu.vn/tin-tuc2", "text": "(*@Thông@*) tin (*@tuyển@*) (*@dụng@*)"}]}
\end{lstlisting}

Các bước xử lý:

\begin{enumerate}
  \item \textbf{Kiểm tra.} \texttt{url} và từng \texttt{out-links.url} phải là url http(s) tuyệt đối, không có thông tin
        đăng nhập; \texttt{format} khớp \texttt{[a-z0-9][a-z0-9.+\_\allowbreak{}-]\{0,31\}}; \texttt{title} không rỗng; phải có
        \texttt{content} hoặc \texttt{base64}, và \texttt{base64} phải giải mã được. Sai thì trả \texttt{400} kèm lý do.
  \item \textbf{Chọn collection theo \texttt{format}.} \texttt{html}/\texttt{htm} vào \texttt{html}; \texttt{jpg},
        \texttt{png}, \texttt{gif}, \texttt{webp}, \texttt{svg}, \texttt{bmp} vào \texttt{picture}; còn lại vào
        \texttt{other}.
  \item \textbf{Upsert theo url} bằng \texttt{findAndModify(upsert=true)}. Crawl lại một trang thì nội dung được ghi đè,
        \texttt{updatedAt} được cập nhật, \texttt{createdAt} giữ nguyên. Hai request cùng url chạy song song có thể cùng
        chèn và gặp \texttt{DuplicateKeyException}; khi đó thao tác được làm lại ở dạng cập nhật.
  \item \textbf{Ghi cạnh đồ thị.} Mỗi liên kết ra là một bản ghi (\texttt{source}, \texttt{destinations}), upsert theo
        cặp nên không có cạnh trùng.
\end{enumerate}

Mất kết nối MongoDB thì trả \texttt{503}, các lỗi ghi khác trả \texttt{500}. Mỗi request được gán một
\texttt{X-Request-ID} và ghi log đầu vào, đầu ra, thời gian (byte Base64 được rút gọn trong log).

\section{Nạp url mầm và theo dõi}

Tệp \texttt{link\_\allowbreak{}all.txt} chứa 101\,461 url mầm thuộc miền \texttt{hust.edu.vn}. Script \texttt{feed.py} chờ
\texttt{mediator} sẵn sàng, gửi danh sách theo lô (mặc định 2\,000 url mỗi lô) rồi gọi \texttt{/start}. Nạp lại nhiều
lần cũng không sao vì url đã có sẽ bị bỏ qua. Sau đó crawl tự mở rộng theo các liên kết mà \texttt{reader} tìm được.

Trang \texttt{dashboard.html} (do dịch vụ \texttt{search} phục vụ tại \path{http://localhost:7700/}) có ba thẻ:
\emph{Theo dõi crawl} (trạng thái \texttt{mediator}, số bản ghi từng collection lấy từ \texttt{/stats}, các bản ghi mới
nhất), \emph{Tìm kiếm} và \emph{Thêm url}. Ở thẻ Thêm url, người dùng dán một hoặc nhiều url; các url được gửi tới
\texttt{mediator} với \texttt{priority: true}, đi qua luồng crawl bình thường, rồi được index ngay bằng
\texttt{POST /index} mà không cần dựng lại cả chỉ mục.

\section{Kết quả thu thập}\label{sec:crawl-stats}

Crawler chạy từ ngày 26/10/2026 đến 31/10/2026. Bảng~\ref{tab:crawl} tổng hợp số liệu từ bản export của MongoDB.

\begin{table}[ht]
\centering
\caption{Số liệu dữ liệu thu thập được}\label{tab:crawl}
\begin{tabularx}{\textwidth}{l r L}
\toprule
Collection & Số bản ghi & Ghi chú \\
\midrule
\texttt{html} & 28\,168 & 95 tên miền con; độ dài nội dung trung vị 1\,079 ký tự, trung bình 2\,575 ký tự \\
\texttt{other} & 4\,690 & 25 tên miền con; PDF 4\,289, DOCX 164, DOC 155, XLSX 61, XLS 16, PPTX 5; độ dài trung vị 4\,795 ký tự \\
\texttt{picture} & 9\,614 & Ảnh kèm metadata và byte Base64 (khoảng 7,5 GB) \\
\texttt{graph} & 84\,673 & 44\,095 nút; 19\,388 trang có liên kết ra; 79,8\% cạnh trỏ vào miền \texttt{hust.edu.vn} \\
\bottomrule
\end{tabularx}
\end{table}

\begin{table}[ht]
\centering
\caption{Các tên miền con có nhiều trang nhất}\label{tab:hosts}
\small
\begin{tabular}{l r @{\hspace{0.6cm}} l r}
\toprule
Host (\texttt{html}) & Số trang & Host (\texttt{other}) & Số tài liệu \\
\midrule
\texttt{bktextile.hust.edu.vn} & 5\,567 & \texttt{hust-dev.hust.edu.vn} & 1\,735 \\
\texttt{clc.hust.edu.vn} & 5\,444 & \texttt{hust.edu.vn} & 1\,418 \\
\texttt{fami.hust.edu.vn} & 3\,806 & \texttt{library.hust.edu.vn} & 680 \\
\texttt{bme.hust.edu.vn} & 1\,535 & \texttt{ctt.hust.edu.vn} & 278 \\
\texttt{icea.hust.edu.vn} & 1\,485 & \texttt{fami.hust.edu.vn} & 205 \\
\texttt{inest.hust.edu.vn} & 1\,396 & \texttt{sep.hust.edu.vn} & 146 \\
\texttt{hust-dev.hust.edu.vn} & 1\,233 & \texttt{sdh.hust.edu.vn} & 67 \\
\texttt{hust.edu.vn} & 981 & \texttt{soict.hust.edu.vn} & 53 \\
\bottomrule
\end{tabular}
\end{table}

Số trang giữa các host lệch nhiều (Bảng~\ref{tab:hosts}): ba host lớn nhất chiếm 53\% số trang web. Nguyên nhân là mỗi
host được thu thập với cùng một tốc độ tối đa, nên host nào có nhiều trang và liên kết dày (website khoa có nhiều bài
tin tức) thì tích lũy được nhiều trang hơn trong cùng thời gian. Hơn một nửa số tài liệu đến từ cổng chính và bản thử
nghiệm \texttt{hust-dev}, nơi đăng các văn bản quy chế, thông báo. Đồ thị có 44\,095 nút nhưng chỉ 19\,388 trang có liên
kết ra; phần còn lại là url được trỏ tới nhưng chưa crawl xong, hoặc là tệp, ảnh, trang ngoài trường.

% ======================================================================
\chapter{Máy tìm kiếm}\label{ch:search}

\section{Tổng quan}

Dịch vụ \texttt{search} (Java 21, Lucene 10.5~\cite{lucene}) có hai quy trình (Hình~\ref{fig:search}): \emph{dựng chỉ mục} từ MongoDB
(lúc khởi động nếu chưa có chỉ mục, hoặc khi gọi \texttt{POST /reindex}) và \emph{trả lời truy vấn} qua
\texttt{GET /search}.

\begin{figure}[ht]
\centering
\begin{tikzpicture}[
  box/.style={draw, rounded corners=2pt, align=left, font=\small, text width=8cm, inner sep=5pt},
  arr/.style={-{Stealth[length=2mm]}, thick}
]
\node[draw, cylinder, shape border rotate=90, aspect=0.25, fill=yellow!15, font=\small, align=center] (m) {MongoDB\\\scriptsize html, other, picture, graph};
\node[box, fill=blue!5, below=0.6cm of m] (b) {\textbf{IndexBuilder.build}\\
1. PageRank trên \texttt{graph}\\
2. Độ hiếm + ngưỡng thung lũng $\rightarrow$ trang rác\\
3. SHA-1 nội dung $\rightarrow$ trang trùng\\
4. Ghi tài liệu Lucene nhiều trường};
\node[draw, cylinder, shape border rotate=90, aspect=0.25, fill=green!10, font=\small, below=0.6cm of b] (i) {Chỉ mục Lucene};
\node[box, fill=blue!5, below=0.6cm of i] (s) {\textbf{Searcher.search}\\
Bước 1: BM25 nhiều trường, lấy 100 ứng viên\\
Bước 2: xếp hạng lại theo cụm từ và PageRank\\
Cắt trang, tạo snippet};
\node[box, fill=gray!8, below=0.6cm of s] (e) {\textbf{Eval}: P@10, nDCG@10 trên nhãn gán tay};
\draw[arr] (m) -- (b); \draw[arr] (b) -- (i); \draw[arr] (i) -- (s); \draw[arr] (s) -- (e);
\end{tikzpicture}
\caption{Hai quy trình của dịch vụ search}\label{fig:search}
\end{figure}

Nguồn dữ liệu được trừu tượng hoá bằng giao diện \texttt{Source} với hai hàm \texttt{records()} và \texttt{edges()}.
\texttt{MongoSource} đọc ba collection nội dung (chỉ các trường \texttt{url}, \texttt{format}, \texttt{title},
\texttt{content}, \texttt{updatedAt}, không bao giờ lấy \texttt{base64}) và collection \texttt{graph}.
\texttt{JsonFileSource} đọc một tệp JSON dạng mảng \texttt{\{url, title, content, olinks\}}; dịch vụ
\texttt{search-json} (cổng 7701) dùng nguồn này để chạy cùng mã trên một bộ dữ liệu khác. Mỗi bản ghi tương ứng một
tài liệu Lucene.

\section{Phân tích văn bản tiếng Việt}

\subsection{Chuẩn hoá Unicode}

Tiếng Việt có hai cách mã hoá cùng một chữ có dấu: ký tự dựng sẵn (``ệ'' là U+1EC7) và ký tự gốc cộng dấu rời (``e''
+ U+0302 + U+0323). Hai chuỗi trông giống hệt nhau nhưng khác byte, nên chữ ``ệ'' gõ ở máy này có thể không khớp với
``ệ'' lưu ở máy khác. Vì vậy mọi chuỗi đều được chuẩn hoá về NFC (\texttt{Analyzers.nfc}) trước khi index và trước khi
truy vấn.

\subsection{Tách từ: unigram âm tiết cộng bigram}

Tiếng Việt viết các âm tiết cách nhau bằng dấu cách, còn một từ có thể gồm nhiều âm tiết (``học phí'', ``bách khoa'').
Nhóm không dùng bộ tách từ riêng (cần thêm mô hình và thư viện) mà dùng \texttt{StandardTokenizer} của Lucene, tách
theo chuẩn Unicode UAX \#29~\cite{uax29} nên mỗi âm tiết thành một token, rồi thêm bigram âm tiết để bắt các từ ghép. Có bốn
analyzer (Bảng~\ref{tab:analyzers}).

\begin{table}[ht]
\centering
\caption{Các analyzer, ví dụ với ``Bách Khoa Hà Nội''}\label{tab:analyzers}
\begin{tabularx}{\textwidth}{l L L}
\toprule
Analyzer & Chuỗi xử lý & Token \\
\midrule
\texttt{vi} & \texttt{StandardTokenizer} $\rightarrow$ chữ thường & \texttt{bách}, \texttt{khoa}, \texttt{hà}, \texttt{nội} \\
\texttt{viBi} & \texttt{vi} $\rightarrow$ \texttt{ShingleFilter(2,2)}, không giữ unigram & \texttt{bách khoa}, \texttt{khoa hà}, \texttt{hà nội} \\
\texttt{viAscii} & \texttt{vi} $\rightarrow$ \texttt{ASCIIFoldingFilter} (bỏ dấu, \texttt{đ}$\rightarrow$\texttt{d}) & \texttt{bach}, \texttt{khoa}, \texttt{ha}, \texttt{noi} \\
\texttt{urlTokens} & Tách theo mọi ký tự không phải chữ hoặc số $\rightarrow$ chữ thường $\rightarrow$ bỏ dấu & \texttt{.../hoc-phi-654580.html} $\rightarrow$ \texttt{hoc}, \texttt{phi}, \texttt{654580}, \texttt{html} \\
\bottomrule
\end{tabularx}
\end{table}

Như vậy ``học phí'' được biểu diễn bằng hai unigram \texttt{học}, \texttt{phí} và một bigram \texttt{học phí}. Tài liệu
chứa đúng cụm ``học phí'' khớp cả ba token nên điểm cao hơn tài liệu chỉ chứa ``học'' và ``phí'' ở hai chỗ rời nhau.
Bigram cũng sinh ra những cặp không phải từ (\texttt{khoa hà}), nhưng các cặp này hiếm khi trùng với truy vấn nên ít
gây nhiễu.

\subsection{Bản bỏ dấu}

Người dùng hay gõ không dấu (``lich thi''), và url, tên tệp trên website của trường cũng thường không dấu
(\texttt{lich-thi-hk1.pdf}). Vì thế mỗi trường văn bản được index thêm một bản bỏ dấu. Bản này có trọng số thấp hơn bản
có dấu, để khi truy vấn có dấu thì trang khớp đúng dấu đứng trên trang chỉ khớp sau khi bỏ dấu (``hồ'' khác ``hổ'', nhưng
bỏ dấu đều thành ``ho'').

\section{Lược đồ chỉ mục}

\begin{table}[ht]
\centering
\caption{Các trường của tài liệu Lucene}\label{tab:fields}
\begin{tabularx}{\textwidth}{>{\raggedright\arraybackslash}p{3.1cm} l l c L}
\toprule
Trường & Loại & Analyzer & Lưu & Mục đích \\
\midrule
\texttt{url}, \texttt{host}, \texttt{collection}, \texttt{format} & String & --- & có & Định danh, lọc theo host và loại \\
\texttt{title} & Text & \texttt{vi} & có & Khớp tiêu đề \\
\texttt{content} & Text (vị trí, offset) & \texttt{vi} & có & Khớp nội dung, cụm từ, snippet \\
\texttt{title\_\allowbreak{}bi}, \texttt{content\_\allowbreak{}bi} & Text & \texttt{viBi} & không & Khớp cụm hai âm tiết \\
\texttt{title\_\allowbreak{}ascii}, \texttt{content\_\allowbreak{}ascii} & Text & \texttt{viAscii} & không & Khớp truy vấn không dấu \\
\texttt{url\_\allowbreak{}tokens} & Text & \texttt{urlTokens} & không & Khớp từ trong url \\
\texttt{sig:pagerank} & FeatureField & --- & không & Tín hiệu liên kết \\
\texttt{updatedAt} & Stored & --- & có & Trả về cho client \\
\texttt{junk} & String & --- & có & \texttt{spam} hoặc \texttt{dup}, chỉ có khi bị đánh dấu \\
\bottomrule
\end{tabularx}
\end{table}

Trường \texttt{content} được index với \texttt{DOCS\_\allowbreak{}AND\_\allowbreak{}FREQS\_\allowbreak{}AND\_\allowbreak{}POSITIONS\_\allowbreak{}AND\_\allowbreak{}OFFSETS}: vị trí cần cho truy vấn
cụm từ ở bước xếp hạng lại, offset cần cho \texttt{UnifiedHighlighter} tạo snippet mà không phải phân tích lại văn bản.

\section{BM25 theo trường}

Điểm văn bản dùng BM25~\cite{robertson2009}. Với truy vấn $q$ và tài liệu $d$ trên trường $f$:
\begin{gather}
\mathrm{BM25}_f(d, q) = \sum_{t \in q} \mathrm{IDF}(t)\cdot
  \frac{\mathit{tf}_{t,d}}{\mathit{tf}_{t,d} + k_1\left(1 - b + b\,\dfrac{|d|_f}{\mathit{avgdl}_f}\right)}, \\
\mathrm{IDF}(t) = \ln\!\left(1 + \frac{N - n_t + 0{,}5}{n_t + 0{,}5}\right)
\end{gather}
trong đó $\mathit{tf}_{t,d}$ là số lần $t$ xuất hiện trong trường $f$ của $d$, $|d|_f$ là độ dài trường,
$\mathit{avgdl}_f$ là độ dài trung bình của trường, $N$ là số tài liệu và $n_t$ là số tài liệu chứa $t$. (Lucene bỏ thừa
số $k_1 + 1$ ở tử số; điều này không làm đổi thứ hạng.)

Mọi trường dùng $k_1 = 1{,}2$. Tham số chuẩn hoá độ dài $b$ khác nhau theo trường: $b = 0{,}75$ cho tiêu đề và url,
$b = 0{,}5$ cho các trường nội dung. Lý do là kho có cả trang web ngắn (trung vị khoảng 1\,000 ký tự) lẫn PDF dài (trung
bình hơn 26\,000 ký tự). Với $b = 0{,}75$, các văn bản quy chế dài bị phạt quá nặng dù chúng thường chính là kết quả
cần tìm.

\section{PageRank}

\subsection{Mô hình}

PageRank~\cite{page1999} là xác suất dừng của một người duyệt web ngẫu nhiên: ở mỗi bước, với xác suất $d = 0{,}85$ người
đó bấm vào một liên kết bất kỳ trên trang hiện tại, với xác suất $1 - d$ thì nhảy sang một trang bất kỳ.
\begin{equation}
\mathrm{PR}(v) = \frac{1-d}{N} + d\left(\sum_{u \to v} \frac{\mathrm{PR}(u)}{\mathrm{outdeg}(u)} + \frac{D}{N}\right),
\qquad D = \sum_{u:\,\mathrm{outdeg}(u)=0} \mathrm{PR}(u)
\end{equation}
Ở đây $D$ là tổng PageRank của các nút cụt (không có liên kết ra), được chia đều cho mọi nút để tổng PageRank luôn bằng 1.

\subsection{Cài đặt}

Đồ thị đọc từ collection \texttt{graph}. Khi nạp, cạnh thiếu đầu mút và cạnh tự trỏ bị bỏ; cạnh trùng bị bỏ bằng khoá
$(a \ll 32) \mid b$ trên chỉ số nút. Mỗi url là một nút, đánh số theo thứ tự gặp, các cạnh lưu trong hai mảng \texttt{int}.
Giá trị khởi tạo là $1/N$; thuật toán lặp tối đa 50 vòng và dừng sớm khi $\sum_v |\mathrm{PR}_{k+1}(v) - \mathrm{PR}_k(v)| <
10^{-6}$. Mỗi vòng tốn $O(N + E)$ và chạy hoàn toàn trong bộ nhớ; với đồ thị 44 nghìn nút và 85 nghìn cạnh, thời gian
tính nhỏ so với thời gian ghi chỉ mục.

\subsection{Đưa vào điểm xếp hạng}

Lúc index, giá trị lưu vào \texttt{FeatureField} là $s = \max(\mathrm{PR} \cdot N,\ 10^{-4})$. Nhân với $N$ để giá trị
trung bình bằng 1, không phụ thuộc kích thước đồ thị; tài liệu không có trong đồ thị nhận giá trị sàn $10^{-4}$. Lúc xếp
hạng lại, PageRank đóng góp qua hàm bão hoà
\begin{equation}
\mathrm{sat}(s) = w \cdot \frac{s}{s + \mathit{pivot}}, \qquad w = 2,
\end{equation}
trong đó \texttt{pivot} do Lucene ước lượng từ phân bố giá trị trong chỉ mục. Hàm này tăng nhanh ở giá trị thấp rồi
phẳng dần, và bị chặn trên bởi $w$. Nhờ vậy trang chủ có PageRank gấp hàng nghìn lần trung bình cũng không lấn át hoàn
toàn điểm văn bản.

\section{Phát hiện trang rác bằng độ hiếm}

\subsection{Quan sát}

Một số trang html trên website của trường bị chèn nội dung rác (từ khoá quảng cáo, văn bản lạ không liên quan tới
trường). Vốn từ của những trang này gần như không trùng với phần còn lại của kho, trong khi trang thật dùng từ vựng
chung của trường (``đào tạo'', ``sinh viên'', ``thông báo''). Từ đó nhóm định nghĩa \emph{độ hiếm} của một trang là IDF
trung bình của các từ trong trang.

\subsection{Độ đo}

Với trang html $p$ có ít nhất 50 từ, gọi $W_p$ là dãy các từ (chuỗi chữ cái \verb|\p{L}+|, chữ thường, NFC) của tiêu đề
và nội dung, $N$ là số trang html đủ dài, và $\mathit{df}(w)$ là số trang html đủ dài chứa $w$:
\begin{equation}
\mathrm{odd}(p) = \frac{1}{|W_p|}\sum_{w \in W_p} \ln\frac{N}{\mathit{df}(w)} .
\end{equation}
Từ phổ biến cho giá trị gần 0, từ chỉ xuất hiện ở một trang cho giá trị $\ln N$. Trang dưới 50 từ không đủ để ước lượng
nên không bao giờ bị coi là rác; tài liệu và ảnh không xét.

\subsection{Chọn ngưỡng tự động}

Ngưỡng phụ thuộc vào kho dữ liệu nên nhóm không đặt bằng tay. Histogram độ hiếm của kho có dạng hai cụm: cụm lớn bên
trái là trang thật, cụm nhỏ nằm cao hơn là trang rác. Ngưỡng được chọn ở đáy thung lũng giữa hai cụm, theo phương pháp
của Prewitt và Mendelsohn~\cite{prewitt1966} (Thuật toán~\ref{alg:valley}).

\begin{algorithm}[ht]
\caption{Chọn ngưỡng độ hiếm theo thung lũng của histogram}\label{alg:valley}
\begin{algorithmic}[1]
\Require điểm độ hiếm $o_1, \ldots, o_m$; độ rộng bin $\beta = 0{,}1$
\Ensure ngưỡng $\theta$ ($+\infty$ nếu không có cụm rác)
\State $h \gets$ histogram của $o$ với bin rộng $\beta$ bắt đầu từ $\min o$; $\mathit{raw} \gets h$
\For{$r = 1$ \textbf{to} $200$}
  \State $P \gets$ các đỉnh của $h$ \Comment{cao hơn bên trái, không thấp hơn bên phải; dải bằng nhau tính một đỉnh}
  \If{$|P| < 2$} \Return $+\infty$ \EndIf
  \If{$|P| = 2$, gọi hai đỉnh là $a < b$}
    \If{$h[b] > h[a]$} \Return $+\infty$ \Comment{cụm bên phải đông hơn thì không phải rác} \EndIf
    \State $v \gets \arg\min_{a \le i \le b} \mathit{raw}[i]$ \Comment{đáy đo trên histogram gốc}
    \If{$\mathit{raw}[v] \le 0{,}5 \cdot h[b]$} \Return $\min o + (v + 0{,}5)\beta$ \EndIf
  \EndIf
  \State $h[i] \gets (h[i-1] + h[i] + h[i+1]) / 3$ với mọi $i$ \Comment{làm mượt}
\EndFor
\State \Return $+\infty$
\end{algorithmic}
\end{algorithm}

Đáy thung lũng được đo trên histogram gốc vì làm mượt nhiều vòng sẽ lấp mất thung lũng. Điều kiện
$\mathit{raw}[v] \le 0{,}5\,h[b]$ loại các thung lũng nông do nhiễu. Trang có độ hiếm vượt ngưỡng được gắn
\texttt{junk:spam}. Ngưỡng của lần dựng gần nhất được ghi vào log và trả ở \texttt{GET /ping}
(\texttt{stats.oddityThreshold}). Cách làm này có giới hạn: nếu số trang rác quá ít, không đủ tạo thành một đỉnh riêng,
thì chúng không bị phát hiện. Khi đó có thể đổi sang ngưỡng cố định theo phân vị.

\section{Loại trang trùng nội dung}

Trang mirror, bản in, bản tiếng Việt và tiếng Anh giống hệt nhau, hay cùng một PDF đăng ở nhiều nơi đều làm trang kết
quả bị lặp. Mỗi tài liệu (cả ba collection) có nội dung từ 200 ký tự trở lên được gán một khoá
\begin{equation}
\mathrm{key}(p) = \mathrm{SHA1}\big(\mathrm{trim}(\mathrm{collapse\_\allowbreak{}ws}(\mathrm{lower}(\mathrm{NFC}(\mathit{content}))))\big).
\end{equation}
Với mỗi khoá có nhiều url, chỉ giữ \emph{một} url: url có PageRank cao nhất, nếu bằng nhau thì lấy url ngắn hơn; trang đã
bị đánh dấu rác không được làm bản giữ lại. Các url còn lại của nhóm bị gắn \texttt{junk:dup}. Cách này chỉ bắt được
trùng khớp hoàn toàn sau chuẩn hoá, không bắt được trang gần trùng.

\section{Quy trình dựng chỉ mục}

\texttt{IndexBuilder.build} dựng lại toàn bộ chỉ mục (\texttt{OpenMode.CREATE}). Để không phải giữ nội dung cả kho trong
bộ nhớ, nó duyệt nguồn dữ liệu nhiều lượt và giữa các lượt chỉ giữ url cùng các con số:

\begin{enumerate}
  \item Đếm cạnh và tính PageRank.
  \item Lượt 1, trên các trang html đủ dài: đếm $\mathit{df}$ của từng từ.
  \item Lượt 2: tính độ hiếm của từng trang html và khoá SHA-1 của mọi tài liệu; chọn ngưỡng rác; chọn bản giữ lại cho
        mỗi nhóm trùng.
  \item Lượt 3: ghi từng tài liệu Lucene kèm cờ \texttt{junk}.
\end{enumerate}

\texttt{POST /reindex} chạy việc dựng ở luồng nền (trả \texttt{202}, đang dựng thì trả \texttt{409}) rồi gọi
\texttt{SearcherManager.maybeRefresh()}, nên tìm kiếm vẫn phục vụ được trong lúc dựng. \texttt{POST /index?url=...} thêm
hoặc thay một trang đã có trong MongoDB vào chỉ mục hiện tại bằng \texttt{updateDocument}; trang thêm lẻ chưa có
PageRank và chưa được kiểm tra rác cho đến lần dựng lại đầy đủ kế tiếp. Sau mỗi lần dựng, log ghi số tài liệu, số trang
rác và ngưỡng, số trang trùng, số cạnh và tỉ lệ tài liệu có mặt trong đồ thị liên kết.

\section{Truy vấn hai bước}

\begin{center}
\small
$q \rightarrow$ NFC, trim $\rightarrow$ truy vấn BM25 nhiều trường $\rightarrow$ 100 ứng viên $\rightarrow$
xếp hạng lại (profile C, D) $\rightarrow$ cắt trang $\rightarrow$ snippet
\end{center}

Việc chia hai bước tách phần rẻ khỏi phần đắt: BM25 trên chỉ mục ngược tính được cho mọi tài liệu khớp, còn truy vấn
cụm từ (phải đọc vị trí) và các tín hiệu ngoài văn bản chỉ cần tính cho 100 ứng viên tốt nhất.

\subsection{Bước 1: lấy ứng viên bằng BM25 nhiều trường}

Với mỗi trường, \texttt{QueryBuilder} phân tích truy vấn bằng analyzer của trường đó rồi nối các token bằng OR (tài
liệu chỉ cần khớp một token). Điểm bước 1 là tổng có trọng số:
\begin{equation}
S_1(d, q) = \sum_{f \in F} w_f \cdot \mathrm{BM25}_f(d, q).
\end{equation}

\begin{table}[ht]
\centering
\caption{Trọng số trường ở bước lấy ứng viên}\label{tab:weights}
\begin{tabular}{l c c c}
\toprule
Trường & Truy vấn có dấu & Truy vấn không dấu (từ B) & Có từ profile \\
\midrule
\texttt{url\_\allowbreak{}tokens} & 1,5 & 1,5 & A \\
\texttt{title} & 3 & --- & A \\
\texttt{content} & 1 & --- & A \\
\texttt{title\_\allowbreak{}bi} & 3 & --- & B \\
\texttt{content\_\allowbreak{}bi} & 1,5 & --- & B \\
\texttt{title\_\allowbreak{}ascii} & 2 & 3 & B \\
\texttt{content\_\allowbreak{}ascii} & 0,7 & 1 & B \\
\bottomrule
\end{tabular}
\end{table}

Các trọng số ở Bảng~\ref{tab:weights} được chọn như sau. Tiêu đề do người viết chọn để tóm tắt bài nên nặng gấp ba nội
dung. Bigram trong nội dung nặng hơn unigram vì khớp cả cụm hai âm tiết là bằng chứng mạnh hơn khớp rời. Bản bỏ dấu nhẹ
hơn bản có dấu để trang đúng dấu đứng trên.

\paragraph{Truy vấn không dấu.} Từ profile B, nếu truy vấn không có dấu nào (tức bằng chính nó sau khi bỏ dấu, ví dụ
``lich thi'') thì chỉ hỏi các trường bỏ dấu (\texttt{title\_\allowbreak{}ascii} trọng số 3, \texttt{content\_\allowbreak{}ascii} trọng số 1) cộng
với \texttt{url\_\allowbreak{}tokens}. Nếu vẫn hỏi cả các trường có dấu, trang chỉ có tên tệp không dấu sẽ được ưu ái hơn trang
``Lịch thi'' viết đúng dấu, trong khi người dùng gõ không dấu chỉ vì tiện chứ không phải vì muốn tìm chữ không dấu.

\paragraph{Bộ lọc.} Tham số \texttt{type} lọc theo \texttt{collection}, \texttt{host} lọc theo host; cả hai là mệnh đề
\texttt{FILTER} nên không tham gia chấm điểm. Profile D thêm hai mệnh đề \texttt{MUST\_\allowbreak{}NOT} để loại tài liệu có
\texttt{junk:spam} hoặc \texttt{junk:dup}.

\subsection{Bước 2: xếp hạng lại}

\texttt{QueryRescorer} xếp hạng lại 100 tài liệu đầu bằng một truy vấn phụ là OR của ba mệnh đề:
\begin{equation}
S_2(d, q) = S_1(d, q) + 2\cdot \mathrm{Phrase}_{\texttt{title}}(d, q) + 1\cdot \mathrm{Phrase}_{\texttt{content}}(d, q)
  + 2\cdot\frac{s_d}{s_d + \mathit{pivot}},
\end{equation}
trong đó $\mathrm{Phrase}_f$ là điểm truy vấn cụm từ trên trường $f$ với độ lệch (slop) bằng 2: các token của truy vấn
phải xuất hiện gần nhau, đúng thứ tự, cho phép chèn tối đa hai vị trí. Mệnh đề cụm từ bị bỏ khi truy vấn chỉ có một
token. Với truy vấn dài hơn hai âm tiết (ví dụ ``công nghệ thông tin và truyền thông''), mệnh đề này kiểm tra cả cụm chứ không chỉ từng cặp bigram, nên việc khớp một tên đơn vị nhiều âm tiết không phụ thuộc vào chỗ bigram có phải từ thật hay không. Thành phần cuối là PageRank bão hoà.

\subsection{Bốn profile xếp hạng}

Các kỹ thuật được bật dần qua bốn profile (Bảng~\ref{tab:profiles}), để đo đóng góp của từng tầng bằng đánh giá offline
trên cùng một chỉ mục. Giao diện mặc định dùng C; khi đánh dấu ``Lọc rác'' thì dùng D.

\begin{table}[ht]
\centering
\caption{Bốn profile xếp hạng}\label{tab:profiles}
\begin{tabularx}{\textwidth}{c L c c}
\toprule
Profile & Lấy ứng viên & Xếp hạng lại & Lọc rác, trùng \\
\midrule
A & BM25 trên \texttt{url\_\allowbreak{}tokens}, \texttt{title}, \texttt{content} & không & không \\
B & A + bigram + bỏ dấu, nhánh truy vấn không dấu & không & không \\
C (mặc định) & như B & cụm từ + PageRank & không \\
D & như B & như C & có \\
\bottomrule
\end{tabularx}
\end{table}

\subsection{Snippet}

Snippet lấy tối đa hai đoạn của \texttt{content} bằng \texttt{UnifiedHighlighter}, bọc từ khớp trong
\texttt{<b>}\ldots\texttt{</b>} và nối các đoạn bằng dấu ``\ldots''. Với truy vấn không dấu chỉ khớp
\texttt{content\_\allowbreak{}ascii}, highlight chạy trên trường bỏ dấu: \texttt{ASCIIFoldingFilter} giữ nguyên offset nên vị trí
trong bản bỏ dấu trùng với vị trí trong văn bản gốc, và dùng lại được để tô đậm trên nội dung có dấu. Nếu vẫn không có
đoạn nào thì lấy 200 ký tự đầu, đã escape HTML.

\subsection{Ví dụ}

Truy vấn ``học phí'' với profile D:
\begin{enumerate}
  \item Lấy ứng viên bằng truy vấn Lucene dưới đây, loại tài liệu rác và trùng, giữ 100 tài liệu đầu:
\begin{lstlisting}
url_tokens:((*@học@*) (*@phí@*))^1.5  title:((*@học@*) (*@phí@*))^3      content:((*@học@*) (*@phí@*))
title_bi:"(*@học@*) (*@phí@*)"^3      content_bi:"(*@học@*) (*@phí@*)"^1.5
title_ascii:(hoc phi)^2   content_ascii:(hoc phi)^0.7
-junk:spam  -junk:dup
\end{lstlisting}
  \item Xếp hạng lại: trang có tiêu đề chứa đúng cụm ``học phí'' được cộng điểm cụm từ nhân 2; trang được nhiều trang
        quan trọng trỏ tới được cộng điểm PageRank.
  \item Trả 10 kết quả đầu kèm snippet.
\end{enumerate}

\section{API tìm kiếm}

\begin{table}[ht]
\centering
\caption{Các endpoint của dịch vụ search}\label{tab:api}
\begin{tabularx}{\textwidth}{l L}
\toprule
Endpoint & Chức năng \\
\midrule
\texttt{GET /search} & Tìm kiếm. Tham số: \texttt{q}; \texttt{type} (\texttt{html}, \texttt{other}, \texttt{picture}); \texttt{host}; \texttt{profile} (A--D, mặc định C); \texttt{page}, \texttt{size} ($\le 50$, $\mathit{page}\times\mathit{size} \le 100$). \\
\texttt{GET /ping} & Số tài liệu, lần index cuối, thống kê (rác, trùng, ngưỡng độ hiếm, số cạnh). \\
\texttt{POST /reindex} & Dựng lại chỉ mục ở nền. \\
\texttt{POST /index?url=} & Index ngay một trang đã có trong MongoDB (404 nếu chưa có). \\
\texttt{GET /} & Giao diện \texttt{dashboard.html}. \\
\bottomrule
\end{tabularx}
\end{table}

\begin{lstlisting}[caption={Ví dụ truy vấn và kết quả}]
$ curl -G localhost:7700/search --data-urlencode 'q=(*@học@*) (*@phí'@*) -d profile=D -d size=10
{"q": "(*@học@*) (*@phí@*)", "profile": "D", "total": 412, "tookMs": 18, "page": 1,
 "hits": [{"url": "https://hust.edu.vn/vi/sinh-vien/hoc-bong-hoc-phi/hoc-phi-654580.html",
           "title": "(*@HỌC@*) (*@PHÍ@*)", "collection": "html", "format": "html",
           "host": "hust.edu.vn", "score": 9.31,
           "snippet": "... (*@mức@*) <b>(*@học@*)</b> <b>(*@phí@*)</b> ...", "updatedAt": 1761200000000}, ...]}
\end{lstlisting}

Trường \texttt{total} là số tài liệu khớp thật (\texttt{IndexSearcher.count}), không chỉ 100. Tuy vậy chỉ 100 kết quả
đầu được xếp hạng nên phân trang dừng ở đó.

% ======================================================================
\chapter{Đánh giá}\label{ch:eval}

\section{Mục tiêu và dữ liệu}

Đánh giá offline so sánh bốn profile A--D trên cùng một tập truy vấn để xem bigram và bỏ dấu, xếp hạng lại, lọc rác có
thật sự cải thiện thứ hạng hay không. Công cụ là lớp \texttt{Eval}, chạy trên cùng thư mục chỉ mục với dịch vụ.

Tập truy vấn \texttt{search/eval/queries.tsv} gồm 25 truy vấn mô phỏng nhu cầu của sinh viên và người ngoài trường, trộn
nhiều dạng:
\begin{itemize}
  \item có dấu, nhiều âm tiết: ``học bổng khuyến khích học tập'', ``chuẩn đầu ra tiếng anh'', ``bảo hiểm y tế sinh viên'';
  \item không dấu: ``lich thi'', ``thong bao nghi tet'', ``thu vien sach dien tu'';
  \item tên đơn vị: ``soict tuyển sinh'', ``viện điện'', ``trường cơ khí'', ``dệt may'';
  \item có năm hoặc tiếng Anh: ``học phí 2025'', ``hust admission''.
\end{itemize}

Tệp nhãn \texttt{qrels.tsv} có các cột \texttt{qid}, \texttt{grade}, \texttt{url}, \texttt{title}, với
\texttt{grade} $\in \{0, 1, 2\}$ lần lượt là không liên quan, liên quan một phần và liên quan hoàn toàn.

\section{Quy trình}

\begin{enumerate}
  \item \textbf{Pooling.} Chạy cả bốn profile cho mỗi truy vấn rồi gộp top 10 của từng profile vào \texttt{qrels.tsv}
        (\texttt{Eval --pool}). Nhãn cũ được giữ, chỉ thêm các cặp (\texttt{qid}, \texttt{url}) mới, nên chạy lại sau khi
        sửa thuật toán không mất công gán nhãn. Pooling~\cite{manning2008} tránh thiên vị về một hệ: tài liệu nào được bất kỳ profile nào
        xếp cao đều được người gán nhãn xem.
  \item \textbf{Gán nhãn tay} vào cột \texttt{grade}.
  \item \textbf{Chấm điểm} bằng \texttt{Eval} (thêm \texttt{--detail} để in điểm từng truy vấn).
\end{enumerate}

\section{Độ đo}

Mỗi profile chạy từng truy vấn không lọc \texttt{type}/\texttt{host} và lấy top $K = 10$. Kết quả chưa có nhãn được tính
là grade 0 và được đếm riêng, để biết số liệu còn thiếu nhãn đến đâu.

\paragraph{Precision@10.} Tỉ lệ kết quả liên quan trong top 10, không xét thứ tự:
\begin{equation}
\mathrm{P@10} = \frac{|\{i \le 10 : g_i \ge 1\}|}{10}.
\end{equation}

\paragraph{nDCG@10}~\cite{jarvelin2002}. Độ lợi là $2^{g} - 1$ (grade 2 có độ lợi 3, gấp ba grade 1), chiết khấu theo vị
trí:
\begin{equation}
\mathrm{DCG@10} = \sum_{i=1}^{10} \frac{2^{g_i} - 1}{\log_2(i + 1)}, \qquad
\mathrm{nDCG@10} = \frac{\mathrm{DCG@10}}{\mathrm{IDCG@10}},
\end{equation}
trong đó IDCG là DCG của thứ hạng lý tưởng, lấy từ \emph{mọi} nhãn đã gán của truy vấn (sắp giảm dần, cắt ở 10). Truy
vấn chưa có nhãn nào có IDCG $= 0$ và nDCG $= 0$, nhưng vẫn được tính vào trung bình. Điểm của một profile là trung bình
trên các truy vấn.

\section{Kết quả}

Bước pooling đã sinh 424 cặp (truy vấn, tài liệu) ứng viên trong \texttt{qrels.tsv}, trung bình khoảng 17 tài liệu khác
nhau cho mỗi truy vấn (tối đa 40 nếu bốn profile không có kết quả nào trùng nhau). Con số này cho thấy top 10 của các
profile khác nhau khá nhiều, tức là các kỹ thuật bổ sung thật sự làm thay đổi thứ hạng.

Tại thời điểm viết báo cáo, cột \texttt{grade} chưa được gán nhãn đầy đủ nên nhóm chưa công bố số liệu P@10 và nDCG@10.
Bảng~\ref{tab:eval} là khuôn kết quả do công cụ \texttt{Eval} in ra; chỉ nên tin các con số khi cột ``chưa gán nhãn''
gần bằng 0.

\begin{table}[ht]
\centering
\caption{Kết quả đánh giá offline trên 25 truy vấn (sẽ điền sau khi gán nhãn)}\label{tab:eval}
\begin{tabular}{c c c c}
\toprule
Profile & P@10 & nDCG@10 & Kết quả chưa gán nhãn \\
\midrule
A & -- & -- & -- \\
B & -- & -- & -- \\
C & -- & -- & -- \\
D & -- & -- & -- \\
\bottomrule
\end{tabular}
\end{table}

Cách đọc kết quả khi đã có số liệu:
\begin{itemize}
  \item \textbf{B so với A} cho biết đóng góp của bigram và bỏ dấu. Kỳ vọng rõ nhất ở các truy vấn không dấu
        (``lich thi'', ``thong bao nghi tet''), vì profile A không có nhánh bỏ dấu nên không khớp được với trang có dấu.
  \item \textbf{C so với B} cho biết đóng góp của cụm từ và PageRank. Kỳ vọng ở các truy vấn nhiều âm tiết (``quy chế
        đào tạo'') và truy vấn tìm trang chính của một đơn vị (``viện điện'', ``soict tuyển sinh'').
  \item \textbf{D so với C} cho biết đóng góp của lọc rác và loại trùng. Kỳ vọng ở các truy vấn mà kết quả bị lặp (cùng
        một thông báo đăng ở nhiều website khoa) hoặc lẫn trang rác.
\end{itemize}

% ======================================================================
\chapter{Triển khai và kiểm thử}\label{ch:ops}

\section{Triển khai}

Toàn bộ hệ thống được mô tả trong một tệp \texttt{compose.yaml}. Mỗi dịch vụ có \texttt{Dockerfile} riêng; dữ liệu bền
vững nằm trên ba volume: \texttt{mongo\_\allowbreak{}data} (MongoDB), \texttt{mediator\_\allowbreak{}state} (hàng đợi SQLite và log) và
\texttt{search\_\allowbreak{}index} (chỉ mục Lucene). MongoDB có healthcheck; \texttt{storage} và \texttt{search} chỉ khởi động khi
MongoDB đã sẵn sàng. Lược đồ \texttt{\$jsonSchema} được nạp tự động khi volume MongoDB còn mới.

\begin{lstlisting}[language=bash, caption={Chạy hệ thống}]
docker compose up -d --build            # (*@khởi@*) (*@động@*) (*@toàn@*) (*@bộ@*)
python3 feed.py                         # (*@nạp@*) 101 461 url (*@mầm@*)
curl localhost:5200/status              # theo (*@dõi@*) (*@hàng@*) (*@đợi@*)
curl localhost:2222/stats               # (*@số@*) (*@bản@*) ghi (*@từng@*) collection
curl -X POST localhost:7700/reindex     # (*@dựng@*) (*@lại@*) (*@chỉ@*) (*@mục@*) sau khi crawl
\end{lstlisting}

Muốn đổi số bản \texttt{downloader} thì sửa hai chỗ cho khớp nhau: \texttt{deploy.replicas} trong \texttt{compose.yaml}
và danh sách \texttt{downloaders} trong cấu hình của \texttt{mediator}. Các tham số vận hành của \texttt{mediator}
(khoảng nghỉ theo host 2 giây, hạn ping 2 giây, hạn load 10 giây, thời gian cho node nghỉ 30 giây, tối đa 5 lần thử,
chu kỳ crawl lại) nằm trong một tệp YAML; cấu hình sai thì dịch vụ báo lý do và không khởi động.

\section{Kiểm thử}

Mỗi dịch vụ có bộ kiểm thử tự động riêng (Bảng~\ref{tab:tests}).

\begin{table}[ht]
\centering
\caption{Kiểm thử tự động}\label{tab:tests}
\begin{tabularx}{\textwidth}{l c L}
\toprule
Dịch vụ & Số ca & Nội dung chính \\
\midrule
\texttt{mediator} & 65 & Chuẩn hoá url, phân loại, lập lịch lại, xoay vòng node, hàng đợi SQLite, API; tích hợp với \texttt{downloader} giả qua HTTP; kiểm thử kiến trúc phân lớp (pytest). \\
\texttt{reader} & 28 & Chọn khối nội dung, văn bản liên kết theo câu, đọc tệp bằng Tika, mã hoá url (JUnit 5). \\
\texttt{storage} & 15 & Kiểm tra đầu vào, chọn collection, upsert, thống kê (JUnit 5). \\
\texttt{search} & 22 & Analyzer, PageRank, độ hiếm và ngưỡng thung lũng, loại trùng, profile, API HTTP, nguồn JSON, độ đo đánh giá (JUnit 5). \\
\texttt{downloader} & 1 & Giao thức \texttt{/ping}, \texttt{/load}, trạng thái bận (\texttt{node:test}). \\
\bottomrule
\end{tabularx}
\end{table}

Ngoài ra \texttt{mediator} có một \texttt{downloader} giả (\texttt{demo/fake\_\allowbreak{}fetcher.py}) trả lời đúng giao thức và chỉ
ghi log các url nhận được. Nó dùng để chạy thử \texttt{mediator} trên một máy và làm mã mẫu cho bên viết
\texttt{downloader}.

% ======================================================================
\chapter{Hạn chế, hướng phát triển và kết luận}\label{ch:limit}

\section{Hạn chế}

\begin{itemize}
  \item \textbf{Chưa dùng văn bản liên kết khi xếp hạng.} \texttt{reader} đã bóc được câu chứa liên kết, nhưng collection
        \texttt{graph} chỉ lưu cặp url nên máy tìm kiếm chưa dùng được anchor text, vốn là tín hiệu mạnh cho truy vấn
        tìm trang cụ thể.
  \item \textbf{Chưa tách từ tiếng Việt.} Bigram âm tiết bắt được phần lớn từ ghép hai âm tiết nhưng không xác định được
        ranh giới từ, và cụm đảo thứ tự (``phí học'') không khớp bigram.
  \item \textbf{PageRank coi mọi liên kết như nhau.} Liên kết menu và chân trang lặp ở mọi trang làm tăng PageRank của
        trang gốc các website; đồ thị cũng chỉ gồm các trang đã crawl.
  \item \textbf{Loại trùng chỉ bắt trùng hoàn toàn.} Hai trang khác nhau một dòng ngày đăng thì không bị coi là trùng.
  \item \textbf{Phát hiện rác phụ thuộc vào phân bố.} Rác quá ít để thành một cụm riêng trong histogram thì không bị
        phát hiện.
  \item \textbf{Phải dựng lại toàn bộ chỉ mục.} Hiện chỉ có thể thêm lẻ từng url; trang thêm lẻ chưa có PageRank và chưa
        qua lọc rác cho đến lần dựng lại kế tiếp.
  \item \textbf{Ảnh chỉ tìm được theo tên tệp.} Không có OCR và chưa dùng \texttt{alt} của ảnh làm nội dung.
  \item \textbf{Phát hiện bài mới còn phụ thuộc cấu trúc website.} Cơ chế duyệt trang danh mục cần biết quy tắc đánh số trang của từng website, và sitemap không phải nơi nào cũng đầy đủ; bài mới nằm ngoài cả hai có thể phải chờ đến chu kỳ crawl lại thông thường.
  \item \textbf{Dữ liệu crawl lệch theo host.} Giãn cách theo host làm host nhiều trang chiếm phần lớn dữ liệu; chưa có
        hạn mức số trang theo host hay ưu tiên theo độ quan trọng (ví dụ tính PageRank từng phần ngay lúc crawl).
  \item \textbf{Chạy trên một máy.} Cả hệ thống triển khai bằng một tệp \texttt{compose.yaml} trên một máy, nên đây là xử lý song song cục bộ chứ chưa phải thu thập phân tán thật sự. Kiến trúc cho phép đặt các bản \texttt{downloader} trên máy khác (chỉ cần khai báo địa chỉ trong cấu hình \texttt{mediator}), nhưng nhóm chưa thử. \texttt{mediator} là một tiến trình duy nhất dùng SQLite nên cũng là điểm nghẽn nếu mở rộng.
  \item \textbf{Chưa kiểm chứng lọc rác và PageRank bằng tay.} Chưa có số liệu về tỉ lệ trang rác bị phát hiện, số trang thật bị loại nhầm, và chưa xem danh sách các trang có PageRank cao nhất để biết có bị chi phối bởi trang điều hướng hay không.
  \item \textbf{Đánh giá chưa xong.} Nhãn liên quan chưa gán đủ nên chưa có số liệu định lượng.
\end{itemize}

\section{Hướng phát triển}

\begin{itemize}
  \item Lưu văn bản liên kết vào \texttt{graph} và index thành trường \texttt{anchor} riêng cho tài liệu đích.
  \item Tách từ tiếng Việt (pyvi, VnCoreNLP) cho trường nội dung, giữ bigram làm trường phụ.
  \item Phát hiện trang gần trùng bằng SimHash hoặc MinHash; giảm trọng số các liên kết điều hướng lặp khi tính PageRank.
  \item Cập nhật chỉ mục từng phần theo \texttt{updatedAt}, tính lại PageRank định kỳ.
  \item Gợi ý khi gõ và sửa chính tả, đặc biệt là khôi phục dấu cho truy vấn không dấu.
  \item OCR cho ảnh và PDF dạng ảnh quét.
  \item Gán nhãn xong, công bố P@10 và nDCG@10, rồi dùng chúng để chỉnh trọng số trường, tham số $b$ và
        \texttt{PR\_\allowbreak{}BOOST}.
\end{itemize}

\section{Kết luận}

Nhóm đã xây dựng một máy tìm kiếm hoàn chỉnh cho miền \texttt{hust.edu.vn}, từ thu thập đến xếp hạng và đánh giá. Bộ
crawler gồm năm dịch vụ đã thu thập hơn 28 nghìn trang web, gần 4\,700 tài liệu và 9\,600 ảnh từ 95 website của trường,
với hàng đợi bền vững, khoảng nghỉ theo host, crawl lại định kỳ, phát hiện bài mới qua sitemap và trang danh mục, và tự tiếp tục khi một khâu gặp lỗi. Máy tìm kiếm xử lý
được các đặc thù của tiếng Việt (chuẩn hoá Unicode, cụm hai âm tiết, truy vấn không dấu), kết hợp BM25 nhiều trường với
cụm từ và PageRank trong mô hình hai bước, và có cơ chế lọc rác, loại trùng không cần đặt ngưỡng bằng tay. Bốn profile
xếp hạng cùng công cụ pooling và nDCG@10 là nền để đo và cải tiến chất lượng xếp hạng một cách có căn cứ khi việc gán
nhãn hoàn tất.

% ======================================================================
\begin{thebibliography}{9}
\addcontentsline{toc}{chapter}{Tài liệu tham khảo}

\bibitem{manning2008}
C.~D. Manning, P.~Raghavan, H.~Schütze.
\newblock \emph{Introduction to Information Retrieval}.
\newblock Cambridge University Press, 2008.

\bibitem{robertson2009}
S.~Robertson, H.~Zaragoza.
\newblock The Probabilistic Relevance Framework: BM25 and Beyond.
\newblock \emph{Foundations and Trends in Information Retrieval}, 3(4):333--389, 2009.

\bibitem{page1999}
L.~Page, S.~Brin, R.~Motwani, T.~Winograd.
\newblock The PageRank Citation Ranking: Bringing Order to the Web.
\newblock Technical report, Stanford InfoLab, 1999.

\bibitem{prewitt1966}
J.~M.~S. Prewitt, M.~L. Mendelsohn.
\newblock The Analysis of Cell Images.
\newblock \emph{Annals of the New York Academy of Sciences}, 128(3):1035--1053, 1966.

\bibitem{jarvelin2002}
K.~Järvelin, J.~Kekäläinen.
\newblock Cumulated Gain-Based Evaluation of IR Techniques.
\newblock \emph{ACM Transactions on Information Systems}, 20(4):422--446, 2002.

\bibitem{lucene}
Apache Software Foundation.
\newblock Apache Lucene 10. \url{https://lucene.apache.org/}

\bibitem{tika}
Apache Software Foundation.
\newblock Apache Tika. \url{https://tika.apache.org/}

\bibitem{uax29}
Unicode Consortium.
\newblock Unicode Standard Annex \#29: Unicode Text Segmentation. \url{https://unicode.org/reports/tr29/}

\bibitem{puppeteer}
Google.
\newblock Puppeteer: Headless Chrome Node.js API. \url{https://pptr.dev/}

\end{thebibliography}

\end{document}
